\section{Framework Architecture}
\label{sec:methodology}

\toolname{} wraps existing RAG components without modifying their internals. The pipeline proceeds through six stages---configuration, data loading, chunking, retrieval, generation, and evaluation---each implemented as an independent module.

\subsection{Configuration and Data Loading}

The framework reads all parameters from a single \texttt{.env} file. Users specify comma-separated lists for each sweepable parameter (LLM models, embedding models, chunk sizes, chunking strategies, retrieval methods, rerankers, prompt templates) and the framework computes their cartesian product. Model identifiers use \texttt{provider:name} syntax (e.g., \texttt{ollama:qwen3:32b}, \texttt{openai:Qwen/Qwen3-32B-AWQ}) with automatic provider detection. Separate base URLs can be configured per role (generation LLM, critic LLM, embedding), enabling distributed setups.

A registry-based adapter pattern maps dataset identifiers to loading functions. Each adapter converts a HuggingFace dataset into a unified \texttt{\{question, ground\_truth, context\}} format. Built-in adapters cover SQuAD~\cite{rajpurkar2016squad}, RagBench, T2-RagBench, Ragas-WikiQA, and RagPerf-Wikipedia-NQ. For datasets with a shared corpus, documents are loaded once and reused across configurations.

\subsection{Chunking}

The framework supports three chunking strategies via LangChain text splitters~\cite{wolf2020transformers}: \emph{recursive} splitting (hierarchy of separators), \emph{fixed-size} splitting (character boundaries with overlap), and \emph{semantic} splitting (groups sentences by embedding similarity at configurable breakpoints). For semantic chunking, chunk size and overlap are ignored since boundaries are determined by embedding similarity.

\subsection{Retrieval and Indexing}

Chunks are embedded and stored in ChromaDB or LanceDB. Vector stores are cached using content fingerprint keys, allowing index reuse across configurations that share the same chunking and embedding setup.

Retrieval supports \emph{similarity search} ($k$ nearest neighbors) and \emph{Maximal Marginal Relevance} (MMR)~\cite{spanakis2018similarity}, which balances relevance with diversity through a tunable $\lambda$ parameter. An optional HyDE step generates a hypothetical answer from the query and uses it as the retrieval input. An optional cross-encoder reranker re-scores retrieved documents and selects the top-$n$ for generation.

\subsection{Generation and Evaluation}

The framework routes requests to Ollama or OpenAI-compatible endpoints based on the model prefix. A wrapper normalizes response formats for RAGAS compatibility. Answer post-processing strips thinking tags from reasoning models.

Each configuration is evaluated through three complementary metric layers:

\textbf{LLM-based quality metrics (RAGAS).} A separate critic LLM assesses faithfulness, context recall, and semantic similarity between generated and ground-truth answers.

\textbf{Information retrieval metrics.} Hit@k, nDCG@k, and Recall@k measure retrieval effectiveness. Context relevance computes cosine similarity between question and chunk embeddings.

\textbf{NLG metrics.} BERTScore~\cite{zhang2020bertscore} (RoBERTa-large), METEOR~\cite{banerjee2005meteor}, ROUGE-L~\cite{lin2004rouge}, and BLEU~\cite{papineni2002bleu} compare generated answers to references. Time-to-first-token (TTFT) and tokens per second (TPS) are recorded per question.

All parameters, metrics, and artifacts are logged to MLflow with a reproducibility bundle capturing the full environment state.
